\newpage
\addcontentsline{toc}{chapter}{Abstract}
\chapter*{Abstract\markboth{Abstract}{}}

Academic timetable preparation is one of the most critical and recurring administrative tasks in every higher education institution. It requires the coordinated allocation of subjects, faculty members, classrooms, laboratories, and student divisions across a fixed weekly schedule while respecting a wide range of constraints such as teacher availability, room capacity, course type (theory or laboratory), and institutional timings. As a college grows in the number of departments, divisions, and courses, the manual preparation of timetables becomes increasingly time-consuming, error-prone, and difficult to maintain. Common issues observed in manual or spreadsheet-based approaches include overlapping faculty assignments, double-booking of rooms, mismatched room types for laboratory sessions, imbalanced faculty workloads, and lack of real-time visibility of schedule updates.

The project titled \textbf{``SamaySetu --- College Timetable Management System''} presents a full-stack, enterprise-grade web application that addresses these challenges through a structured, constraint-aware, and role-based scheduling platform. Rather than attempting opaque algorithmic auto-generation, the system empowers the academic coordinator with an intuitive visual timetable builder combined with a robust eight-point conflict detection engine. This hybrid approach preserves the coordinator's domain expertise while eliminating the mechanical errors inherent to manual scheduling.

The backend is developed using \textbf{Java 17 with Spring Boot 3.5.5} and secured with \textbf{JWT-based authentication, BCrypt password hashing, Spring Security role-based authorization, rate limiting, and global input sanitization}. All academic data is persisted in a \textbf{PostgreSQL 17} relational database, and frequently-read published timetables are cached in \textbf{Redis} for low-latency access. The frontend is implemented in \textbf{React 18 with TypeScript, Vite, Tailwind CSS, Zustand for state management}, and \textbf{dnd-kit for accessible drag-and-drop} interaction. Timetable exports are generated on the server using \textbf{OpenPDF} for PDF and \textbf{Apache POI} for Microsoft Excel.

The system supports four distinct user roles --- Administrator, Head of Department, Timetable Coordinator, and Teacher --- each with a scoped set of permissions enforced at both the URL and method level. The timetable builder supports manual entry placement, drag-and-drop movement and swapping, a dedicated Lab Session Wizard that atomically creates parallel batch entries across two consecutive periods, a semester filter, a copy-from-division feature, and a pre-publish validation dashboard that reports blocking errors and informational warnings before any draft is made live. Published timetables are cached in Redis and invalidated automatically when an administrator re-publishes. Comprehensive audit logging, security headers, and an environment-driven configuration model make the system production-ready for deployment on AWS EC2 with AWS Amplify as the frontend host.

Overall, the \textbf{SamaySetu} platform significantly reduces the manual effort involved in timetable preparation, eliminates the most common classes of scheduling errors through automated conflict prevention, improves transparency for faculty and students through real-time access, and establishes a scalable architectural foundation on which future enhancements such as student-facing portals, examination scheduling, and constraint-solver-based auto-generation can be built.


\vspace{10mm}
\hrule
